\documentclass[11pt]{article}
\usepackage[margin=1in]{geometry}
\usepackage{amsmath,amssymb}
\usepackage{hyperref}
\usepackage[numbers,sort&compress]{natbib}

\newcommand{\code}[1]{\texttt{#1}}

\title{Hyperframes and the Epistemology of Machine-Gathered Semantics\\
\large Coherence, not truth, and where that leads mathematically}
\author{Project \texttt{thejeb}}
\date{\today}

\begin{document}
\maketitle

\begin{abstract}
Set theoretic estimation (STE) as mechanized in this project does not
start from a crisp assertion ``$X = v$.''  It starts from a
\emph{possibility set}: a source commits only to ``the referent is
somewhere in $S \subseteq \Xi$.''  We call this primitive object a
\emph{hyperframe}, and argue it is the right primitive precisely
because the alternative --- pretending a text (or a labeler reading a
text) hands us a single determinate referent --- is false to how
meaning is actually gathered, by machine or by hand.  This note states
that epistemological position plainly, brackets two rival pictures of
meaning it is \emph{not} chasing (a Platonic universal meaning, a
Chomskyan deep structure resolved in the mind), and is honest about
where the project's formal guarantees stop and interpretation begins.
It then works out the mathematics that dropping ``truth as
correspondence'' and keeping ``coherence as criterion'' actually
forces: lattice theory and Formal Concept Analysis, rough-set
indiscernibility, monotone update semantics, a paraconsistent reading
of local contradiction, and a genuine metric geometry on readings.
Every claim marked as a \emph{theorem} is machine-checked in the
project's Lean library and cited by name; everything else is framing,
cited to the literature it borrows from.
\end{abstract}

\section{The primitive is the hyperframe, not the crisp frame}
\label{sec:primitive}

A \emph{crisp frame} says a variable $X$ takes a specific value $v$.
That is the textbook picture, and it is the wrong starting point for
text.  When a document is read --- by a person or by a language model
standing in for one --- what is actually extracted is rarely ``$X =
v$.''  It is closer to ``on the evidence in this passage, $X$ is one
of $S$,'' where $S$ may be a singleton, but need not be.  The reader
does not know, and the text does not always determine, the exact
referent; reference is \emph{underdetermined by design}, not by
accident of a lazy extractor. This is Quine's inscrutability of
reference in a concrete engineering guise: even fully cooperative,
fully competent interpretation leaves room for more than one referent
consistent with everything actually said \citep{quine1960word}.

The project's Lean development takes this underdetermination as the
object level, not as noise to be denoised away. In
\code{Ste.HyperFrame}, a source contributes a \emph{possibility set}
\code{S i \ensuremath{\subseteq} \ensuremath{\Xi}} over a space of
candidate referents \code{\ensuremath{\Xi}}, and satisfaction is
literally membership:
\[
  \code{sat S a i := a \ensuremath{\in} S i}.
\]
A crisp frame is the special case $|S\,i| = 1$; nothing downstream
requires that specialization. Estimation over a corpus is then the
question this note is really about: do the possibility sets contributed
by many sources cohere, and what does the object built from that
coherence look like?

\section{What is bracketed, and why}
\label{sec:bracket}

Two rival pictures of ``the meaning'' are conspicuously absent from
this project, and their absence is a decision, not an oversight.

\textbf{Not a Platonic meaning.} One may privately hold that beneath
any author's words there is a mind-independent fact of the matter ---
an ideal Tree, an ideal Justice --- that the text either does or does
not correctly describe. This note takes no position on that question
and needs none. It is bracketed, in the sense Husserl gave that word:
the metaphysical question is suspended, not answered, because the
object under study is not the world the text is about but the meaning
the text \emph{commits to} \citep{husserl1913ideas}. Call this a
methodological \emph{epoch\'e}: whatever one believes about
correspondence to an external reality, it plays no role in what STE
checks.

\textbf{Not a Chomskyan universal deep structure either.} A different,
non-Platonic universalism is also declined: the idea that every
utterance resolves, in a competent mind, to one shared underlying
representation across speakers, languages, and eras --- a deep
structure of which surface texts are mere transformations
\citep{chomsky1965aspects}. That picture chases a single semantics
\emph{behind} the corpus. This project chases the opposite: the
semantics \emph{of} the corpus, as specified, document by document, by
what each document's author actually committed to. Two corpora need
not, and generally will not, resolve to the same underlying meaning
even for words that look alike on the page, because there is no
``underneath'' being posited at all.

What is left, once both universalisms are set aside, is what the
Cambridge School of intellectual history calls the object of the
exercise: \emph{situated authorial commitment} --- what a given writer,
in a given idiom, at a given moment, was doing with these words, which
is recoverable, if at all, only from context: the text, its
contemporaries, its conventions \citep{skinner1969meaning}. Recovering
it is an act of interpretation that only makes sense from inside a
horizon of prior understanding brought to the text and revised against
it --- the Gadamerian hermeneutic circle
\citep{gadamer1960truth}. Under this scoping, truth-as-correspondence
to an external world is simply the wrong target: a corpus of internally
coherent falsehoods (a consistent conspiracy theory, a self-contained
myth cycle, a novel's invented history) is a perfectly good object of
study, and STE should say so --- \emph{coherent}, not \emph{true} ---
without complaint. That is a feature of the scoping, not an
embarrassment to be argued around.

\section{The extractor is a fallible, situated labeler}
\label{sec:extractor}

Nothing about using a large language model to pull frames out of a
document changes this picture; it only makes the mechanism explicit.
An LLM extractor stands in for a human labeler doing the same job by
hand, and a human labeler is exactly as fallible and exactly as
situated as the model: both bring priors, both make judgment calls
about what a passage commits its author to, and neither has privileged
access to an author's mind.

Concretely: read Spenser's \emph{The Faerie Queene} against Rowling's
\emph{Harry Potter} and try to extract a shared frame schema for, say,
\textsc{quest} or \textsc{wizard}. The two texts will produce
\emph{materially different} frames --- different slot structures,
different obligatory fillers, different implicit background
assumptions --- not because one extraction is more careful than the
other, but because the frames genuinely are specified by, and bounded
by, each text's own idiom, allegorical machinery, and narrative
conventions. There is no neutral, text-independent \textsc{wizard}
frame that both authors are imperfectly instantiating. The object of
study, stated without hedging, is: \emph{context-bounded texts, and the
machine (or human) gathering of the semantics as specified by those
texts.} Nothing more universal is on offer, and nothing more universal
is being claimed.

\section{Two levels of coherence, and only one is machine-checked}
\label{sec:twolevel}

Once frames are extracted, two different questions can be asked about
them, and the project is careful to keep them separate.

\begin{description}
\item[Fidelity (frame $\leftrightarrow$ document).] Does an extracted
  possibility set $S_i$ faithfully render what document $i$ actually
  commits to? This is an act of interpretation. It is not, and in
  principle cannot be, kernel-checkable: no proof assistant can compare
  a formal set to a paragraph of English (or Elizabethan English) and
  certify that the set says what the paragraph means. The extractor,
  human or machine, is fallible, and an extracted frame \emph{might not}
  cohere with its document even when every downstream computation is
  flawless.
\item[Consistency (frame $\leftrightarrow$ frame).] Do the extracted
  frames cohere with one another --- is $\bigcap_i S_i \ne
  \varnothing$, or, more generally, is the corpus feasible on some
  constraint set? This is exactly STE feasibility, and it \emph{is}
  machine-checked: it is a computation over sets, and the project's
  Lean development proves theorems about it.
\end{description}

The honest statement of what the machinery guarantees is therefore a
conditional, not an unconditional truth claim:
\begin{quote}
\emph{If} the extracted frames faithfully render what their documents
commit to, \emph{then} the feasibility verdict STE returns is sound.
\end{quote}
The consequent is a theorem. The antecedent is an empirical, auditable,
non-formal claim about extraction quality that the project does not
and cannot discharge inside the kernel. This is worth quarantining in
plain sight, the same way the project already keeps
\code{native\_decide} out of its trusted proof base: not because the
excluded step is worthless, but because pretending it is
kernel-strength would be dishonest about what has actually been
checked.

This scoping is productive, not merely modest. Because fidelity and
consistency are separated, STE can partially \emph{audit} its own
extractor rather than simply trust it. Two independent renderings of
the \emph{same} document that STE-conflict are pure extraction noise
--- a fidelity failure, localized and re-checkable, and clearly
distinguishable from genuine cross-\emph{document} authorial
disagreement (which is consistency information, not noise). And
because frames can be required to carry provenance spans back to the
source text, frame$\leftrightarrow$document coherence is made
\emph{inspectable} by a human reviewer even though it is not, and will
not be, formally provable.

\section{Where that leads us, in a math kind of way}
\label{sec:math}

Bracketing truth-as-correspondence and adopting coherence-as-criterion
is not a retreat into vagueness. It relocates the target of formal
work: meaning is no longer analyzed as a proposition that is true or
false of the world, so the natural home for it is no longer the logic
of propositions. It becomes an object in \emph{order theory} and
\emph{closure systems} instead --- and that relocation has real,
provable content. Five consequences follow, each with a machine-checked
piece already in the Lean development.

\paragraph{1. Meaning is a lattice object (Formal Concept Analysis).}
The STE satisfaction relation \code{sat S a i := a \ensuremath{\in} S i} is, on
the nose, a \emph{formal context} in the sense of Ganter and Wille
\citep{ganter1999fca}: hypotheses ($a \in \Xi$) as objects, constraints
($i \in I$) as attributes. \code{Ste.HyperFrame} proves that the
partial feasibility set on a constraint window $J$ is
\emph{definitionally} the lower polar of that context
(\code{partialFeasibilitySet\_eq\_lowerPolar}), so the hyperframes of a
corpus --- pairs of a feasible-hypothesis set and its matching
constraint signature --- are exactly the \emph{formal concepts} of the
context and form a complete lattice (the abbreviation \code{HyperFrame}
for Mathlib's \code{Order.Concept}, with its inherited
\code{CompleteLattice} instance). Every partial feasibility set is a
concept extent (\code{isExtent\_partialFeasibilitySet}), and every
concept's extent recovers the feasible set of its own intent
(\code{hyperFrame\_extent\_eq\_feasibilitySet}). Meaning, on this
picture, is not a truth-value assignment; it is a point in a concept
lattice, and STE feasibility \emph{is} Formal Concept Analysis.

\paragraph{2. Reference is a rough-set indiscernibility floor.}
Two candidate referents can be indistinguishable by every constraint a
corpus happens to supply, without being the same referent. \code{Indisc
S a b := \ensuremath{\forall} i, a \ensuremath{\in} S i \ensuremath{\leftrightarrow} b
\ensuremath{\in} S i} captures this: it is an equivalence relation
(\code{indisc\_equivalence}), and it is exactly Pawlak's
indiscernibility relation for the STE formal context
\citep{pawlak1982rough}. No subfamily of constraints can ever separate
indiscernible hypotheses (\code{indisc\_mem\_iff}) --- this is the
resolution floor of the corpus, below which reference is irreducibly
underdetermined, not a defect to be engineered away. Reference, on this
picture, is not delivered as a single object; it is a fixed point of
constraint accumulation, and the indiscernibility classes are exactly
the fibers of the constraint-signature map on the concept lattice
(\code{indisc\_iff\_upperPolar\_singleton\_eq}).

\paragraph{3. Content is monotone and dynamic.}
Because $\bigcap_{i \in K} S_i \subseteq \bigcap_{i \in J} S_i$
whenever $J \subseteq K$, pure accumulation of evidence can only shrink
what remains feasible; over a finite hypothesis space this is
cardinality-monotone
(\code{ncard\_partialFeasibilitySet\_antitone}), and it disproves a
naive worry about oscillation: if enforcing a larger constraint set
somehow yields a \emph{larger} feasible set, that set was not actually
larger --- some constraint must have been \emph{dropped}
(\code{lt\_ncard\_imp\_not\_subset}). Growth in what is deemed feasible
therefore always signals retraction, never pure accumulation, which is
exactly the shape update semantics gives information states: possibility
sets that only ever shrink on genuine new information and require
explicit retraction otherwise \citep{veltman1996defaults}. Enforcing two
constraint sets together is enforcing their union, and feasibility sets
compose by intersection under that union
(\code{partialFeasibilitySet\_union}) --- readings superpose rather
than overwrite.

\paragraph{4. Contradiction is local and paraconsistent, not
explosive.}
A possibility set $S_i(v)$ for a variable $v$ can be a singleton
(determinate), empty (over-determined --- a local contradiction), or
have more than one element (under-determined). This is exactly Belnap's
four-valued generalization of classical truth
\citep{belnap1977computer}, transported from the two-element set
$\{\text{true}, \text{false}\}$ to an arbitrary reference space $\Xi$:
none, one, or many admissible values, instead of just true or false.
Crucially, emptiness on one variable does not trivialize the rest of
the corpus --- there is no principle of explosion here, in the spirit
of da Costa's paraconsistent logics \citep{dacosta1974paraconsistent}.
When the \emph{global} intersection across all sources is empty, the
right object to hand back is not silence but the concept lattice of
maximal jointly-feasible sub-readings --- an admissible-precisification
move in the sense of van Fraassen's supervaluationism
\citep{vanfraassen1966supervaluations}: keep every way of resolving the
conflict that is itself internally coherent, and let disagreement show
up as multiplicity in the lattice rather than as global collapse.

\paragraph{5. Content has geometry.}
The lattice of readings is not just ordered; it carries a genuine
metric. \code{Ste.PartitionRank} equips the coreference partition
lattice with a combinatorial, entropy-free rank
$r(P) = |\alpha| - {}$blocks$(P)$ and proves it \emph{submodular}
(\code{rank\_submodular}), the counting shadow of Shannon's 1953
observation that entropy is a submodular rank function on the lattice
of partitions \citep{delsol2024}; this rank is monotone in the corpus
(\code{mustSetoid\_eq\_sInf}, \code{rank\_mustSetoid\_mono}), so more
documents can only add forced identifications, never remove them.
\code{Ste.InfoDistance} turns this rank into an honest metric, the
counting shadow of Shannon's information distance: symmetric,
vanishing exactly on equal readings (\code{shannonDist\_eq\_zero\_iff}),
and satisfying the triangle inequality
(\code{shannonDist\_triangle}) --- so ``how far apart are two
readings'' is a well-posed question with a real answer, not a metaphor.
On top of that metric, \code{Ste.CouplingRank} makes ``genuine
cross-topic coupling'' measurable rather than impressionistic: the
excess rank a partition carries once forced through another
(\code{couplingExcess}) is provably positive exactly when the
partitions are not already aligned
(\code{couplingExcess\_pos\_of\_not\_le}). And where the corpus really
is uncoupled --- rectangular, in the project's terms --- feasibility
becomes a purely local check with no obstruction to gluing
(\code{Ste.HellyConsistency}'s \code{frame\_hasHelly\_two}: pairwise
consistency of authors already certifies whole-corpus consistency, a
frame-theoretic Helly number of 2; \code{Ste.CechCover}'s
\code{rectangular\_cechVanishesCover} and \code{Ste.CechObstruction}'s
\code{diagonal\_cechObstruction} for the discriminating example where
that vanishing genuinely fails). The correctness criterion throughout
is \emph{internal} coherence --- a consistency and closure condition
on a lattice --- never external correspondence to a fact of the matter
outside the corpus.

\section{The through-line}
\label{sec:conclusion}

None of this is a way of avoiding hard questions about meaning by
retreating into formalism. It is the opposite move: bracketing
truth-as-correspondence is what \emph{lets} the hard questions be asked
in a setting that has theorems --- lattices, closure systems,
rough-set indiscernibility, paraconsistent possibility spaces, and a
genuine metric geometry on readings --- instead of a setting
(propositional truth of an authorial intention no one can access) that
mostly has aphorisms. Section~\ref{sec:twolevel}'s conditional is the
hinge the whole argument turns on: STE never claims a document's frame
is correct, only that \emph{if} extraction was faithful, feasibility is
sound, and the machinery is built so that claim can be inspected,
partially audited, and never quietly upgraded to more than it is.
Where the fidelity question stops --- at the boundary of what a
possibility set can say about a paragraph of Spenser or a paragraph of
Rowling --- is exactly where situated, competent, fallible
interpretation has to take over, the way it always has.

\bibliographystyle{plainnat}
\bibliography{../../refs}

\end{document}
